\chapter{电子转移推动的反应}

\begin{kaichang}
找一枚生了锈的铁钉——阳台上、旧自行车上、小区健身器材的螺丝上，多半都有。把它和一枚新钉子放在一起掂一掂、看一看：旧的表面裹着一层红棕色的、疏松的壳，新的则是灰亮光滑的金属。

现在请你猜三个问题。第一：铁变成铁锈，是变重了还是变轻了？很多人凭直觉说"烂掉了，当然变轻"。第二：这层红棕色的壳，是铁"变质"出来的，还是从外面"贴上去"的？第三：苹果切开后放一会儿，切面会变成黄褐色；切开的土豆也是。一件发生在金属上、一件发生在水果上，它们之间有没有关系？

先写下你的猜测。这一章结束时你会发现：生锈、苹果变褐、火柴燃烧、漂白剂褪色、你此刻的呼吸——这些看似毫不相干的事，在粒子层面是同一种戏码：\textbf{电子从一个原子手里，转到了另一个原子手里}。
\end{kaichang}

\section{看得见的"变脸"}

先做纯观察，不急着解释。把下面几件事摆在一起看：

\begin{itemize}[itemsep=2pt]
\item 铁钉在潮湿的空气中，几天就冒出红棕色的锈斑；放在干燥处或涂了油，就锈得很慢。
\item 一根火柴划过磷面，"嚓"地燃起明亮的火焰，烧完只剩一小截黑色炭棍，冒出白烟。
\item 切开的苹果，切面朝空气放半小时，从白变成黄褐；泡进盐水或柠檬汁里，变色就慢得多。
\item 沾了蓝黑墨水的白布，滴上稀释的漂白剂，颜色眼看着褪去。
\item 把一根光亮的铜丝插进硝酸银溶液（实验室演示），铜丝表面慢慢"长"出一层银白色的绒毛，溶液则从无色渐渐变蓝。
\end{itemize}

五件事，五种颜色变化，看起来风马牛不相及。但它们有一个共同的嫌疑犯：\textbf{氧或者电子}。铁锈里有氧，燃烧要氧气，苹果变褐要接触空气……那铜丝"长银"呢？溶液里并没有氧气参与。还有火柴烧完——如果称一称，生成的烟和灰加起来比原来的火柴棍\textbf{更重}。多出来的质量从哪来的？

1770 年代的化学家们就卡在这个问题上。卡住他们的，是一个听起来很合理的旧理论，我们留到证据故事里去讲。先自己把镜头推进到粒子层。

\section{镜头推进：是电子换了主人}

回忆第 9 章的图景：原子核外住着电子，最外层电子决定元素的脾气。第 12 章讲过，钠这类金属最外层电子"待得不稳"，容易脱手；第 13 章讲过，氯这类卤素"见电子就抢"。第 18 章则讲过它们交接电子之后的事：钠变成带正电的钠离子，氯变成带负电的氯离子，靠静电吸引抱成食盐。

现在把这个画面搬到燃烧的镁条上。镁在空气中剧烈燃烧，发出耀眼的白光（烟花和闪光灯里就有它），生成白色粉末氧化镁。粒子层面发生了两件事：每个镁原子交出 2 个电子，每个氧原子收下 2 个电子，然后带电的它们抱在一起。镁和氧谁也没有"变质"成别的元素——质子数都没动——只是\textbf{电子换了主人}。

再看铜丝和硝酸银。溶液里的银离子 \chem{Ag^+} 碰到铜原子，会从铜身上"收走"1 个电子，自己变回中性的银原子，沉积在铜丝表面——这就是那层银白色绒毛。而失去电子的铜变成 \chem{Cu^{2+}} 离开表面进入溶液，把溶液染成蓝色。注意：整个过程\textbf{没有氧参加}，可它和镁燃烧是同一类反应——都有电子的转移。

于是化学家给了这类反应一个统一的名字：\textbf{氧化还原反应}。判断标准只有一条：

\begin{itemize}[itemsep=2pt]
\item 一个粒子\textbf{失去}电子，就说它被\textbf{氧化}了；
\item 一个粒子\textbf{得到}电子，就说它被\textbf{还原}了。
\end{itemize}

这个名字起得其实很不幸。"氧化"两个字来自氧气的"氧"，因为氧气是历史上第一个被认出来的"抢电子高手"，于是人们用它的名字命名了整类反应。但名字只是历史的化石。铜被银离子氧化，没有氧；氢气在氯气里安静地燃烧生成氯化氢 \chem{HCl}，没有氧，照样是氧化还原——氢失了电子，氯得了电子。\textbf{氧化还原不等于得氧失氧}，这是本章的第一条红线。判断时要盯住的是电子的去向，不是氧的出没。

还有一条记账纪律：电子不会凭空消失，也不会无家可归地飘在溶液里。一个反应里，某处失去的电子总数，必然等于另一处得到的电子总数。氧化和还原永远是\textbf{成对发生}的，像一枚硬币的两面——所以"氧化"和"还原"从来不能单独说一个物质"被氧化了"就完事，必须同时问：谁被还原了？

由此引出两个常听到的称呼：把电子\textbf{给出去}、让别人被还原的那位，叫\textbf{还原剂}；把电子\textbf{收下来}、让别人被氧化的那位，叫\textbf{氧化剂}。镁燃烧中，镁是还原剂（自己被氧化），氧气是氧化剂（自己被还原）。听着有点绕？记住一个画面就好：\textbf{剂}是"动手的那个人"，它的名字说的是它对\textbf{别人}做的事，而它自己遭遇的恰好相反。氧气天天把别的物质氧化，所以它叫氧化剂。

拿你划过的那根火柴把这两个角色认一遍。火柴头里藏着氧化剂氯酸钾 \chem{KClO_3}（一个囤积了大量氧的"电子买家"），还有还原剂硫和磷的化合物；火柴盒侧面的红磷涂层则是另一份还原剂。划火柴的那一下，摩擦生热提供了第 29 章讲的活化能，把反应推过能量山头——接下来氧化剂和还原剂之间电子大搬家，放出的热自己维持火焰。一个小小的火柴头，就是一个包装好的、等你触发的氧化还原反应。安全火柴把氧化剂和一部分还原剂分放在两头（火柴头和盒边磷面），正是利用"反应需要双方碰面"这条规则来防火——氧化还原的知识，就这样藏在"为什么要划一下才着"这件小事里。

\section{给电子记账：氧化数}

电子转移说起来清楚，真到复杂分子里就麻烦了。高锰酸钾 \chem{KMnO_4} 里的锰、硫酸 \chem{H_2SO_4} 里的硫，电子在共价键里是"共享"的（第 19 章），并没有整颗电子真的搬家，这账怎么算？

化学家的办法是引入一套记账工具：\textbf{氧化数}（也叫氧化态）。规则的核心思想是"假装清算"——把每一根共价键里的共享电子，全部判给电负性更大（抢电子能力更强，见第 11 章）的那个原子，然后数一数每个原子"名下"的电子比中性时多了还是少了几个，多几个记负几，少几个记正几。几条常用记账规则：

\begin{itemize}[itemsep=2pt]
\item 单质里原子的氧化数是 0：\chem{Fe}、\chem{O_2}、\chem{Cu} 都算 0。
\item 单原子离子的氧化数就是它的电荷：\chem{Na^+} 是 $+1$，\chem{Cl^-} 是 $-1$。
\item 化合物里，氧通常算 $-2$（过氧化物等例外），氢通常算 $+1$（金属氢化物例外）。
\item 一个中性分子里所有原子的氧化数加起来等于 0；一个离子里加起来等于它的电荷。
\end{itemize}

拿高锰酸钾算一笔：钾是 $+1$，四个氧各 $-2$ 共 $-8$，总和要为 0，锰只能是 $+7$。再算呼吸作用的主角葡萄糖燃烧：\chem{C_6H_{12}O_6} 里氢 $+1$、氧 $-2$，六个碳加起来必须是 $0-(+12)-(-12)=0$，平均每个碳是 0；烧完变成 \chem{CO_2}，碳变成 $+4$。每个碳"亏"了 4 个电子——葡萄糖被氧化了，这正是你的身体从糖里取能量的方式（第 52、53 章细讲）。

现在说本章的第二条红线：\textbf{氧化数是记账工具，不是贴在原子上、可以直接测量的小标签}。\chem{MnO_4^-} 里的锰并不是真的带着 $+7$ 的电荷——那会是 7 颗电子的彻底搬家，实际上锰和氧之间的电子是共享的，只是分账时多判给了氧。真实的电荷分布要靠量子力学计算和精密的实验（比如 X 射线光谱）去逼近，往往是不漂亮的分数。氧化数的价值不在于"真实"，而在于\textbf{好用}：它让电子账一目了然，谁升谁降、升降多少，一眼看穿，配平方程式时尤其顺手。把它当成复式记账法里的科目编号，而不是保险柜里的现金，你就用对了它。

用氧化数的语言重说前面的定义，就不再依赖"电子真的搬没搬家"：氧化数升高 $=$ 被氧化，氧化数降低 $=$ 被还原。碳从 0 升到 $+4$，锰从 $+7$ 降到 $+2$（高锰酸钾褪色实验里），账就这样记。

\begin{qiaoliang}
一枚铁钉生锈，要转移多少个电子？取一枚约 \ud{5}{g} 的铁钉，铁的摩尔质量约 \ud{56}{g/mol}，所以它含约 $5/56 \approx 0.09\ \mathrm{mol}$ 铁原子。铁锈的主要成分里铁是 $+3$ 价，即每个铁原子交出约 3 个电子。那么整枚钉子锈透，共交出约 $0.27\ \mathrm{mol}$ 电子，乘以 \ud{6.02\times 10^{23}}{mol^{-1}}，约 $1.6\times 10^{23}$ 个电子——一千六百亿亿个电子，从铁原子的手里，一个一个搬到了氧原子的手里。而你家厨房里每小时从煤气（或空气与食物）里搬走的电子，也是这样的天文数字。化学变化听起来温和，账面上的电子流却从来汹涌。
\end{qiaoliang}

\section{拆开两半看：半反应与电子守恒}

账要记平，最好的办法是把氧化和还原\textbf{拆成两个半反应}来写。还看铜丝和硝酸银。铜这边：

\[\chem{Cu} \rightarrow \chem{Cu^{2+}} + \chem{2e^-}\]

银那边：

\[\chem{Ag^+} + \chem{e^-} \rightarrow \chem{Ag}\]

铜每交出 2 个电子，需要 2 个银离子来接。于是把第二个半反应整体乘以 2，再相加，电子正好抵消：

\[\chem{Cu} + \chem{2Ag^+} \rightarrow \chem{Cu^{2+}} + \chem{2Ag}\]

这就是\textbf{半反应法}。它的纪律只有一条：\textbf{电子守恒}——左边交出的电子数，必须等于右边收下的电子数。配上第 26 章学过的原子守恒，再检查电荷两边相等，一个氧化还原方程式就配平了，而且你清楚地知道每一颗电子从谁手里到了谁手里，而不是靠"试数字"碰运气。

半反应法还有个隐藏的好处，先在这里埋一句话：既然两个半反应可以分开写，它们能不能在空间上也分开——让铜在这头失电子，银离子在那头得电子，电子绕一段\textbf{外电路}走过去？能。那就不是烧杯里的反应了，那是电池。第 34 章整章讲这件事。

\section{一场被天平终结的争论}

\begin{zhengju}
18 世纪流行的燃烧理论叫"燃素说"：可燃物里含有一种叫"燃素"的东西，燃烧就是燃素逃进空气。这个理论能解释火焰和烟，看起来挺顺。但它有一个越来越扎眼的麻烦：金属烧过之后剩下的"烧渣"（今天叫氧化物）\textbf{比原来的金属更重}。逃掉了东西，怎么反而变重了？燃素说的支持者只好打补丁：燃素具有"负的重量"。一个理论开始靠"负重量"这种东西续命，离出事就不远了。

1770 年代，法国化学家拉瓦锡把天平请到了舞台中央。他把锡和铅放在密封容器里加热，称总质量——不变；打开容器，空气涌进去，再称——增加的质量恰好等于金属变成烧渣所增加的质量。他又设法把汞的烧渣加热分解，收集到一种气体，质量对得上，而且蜡烛在这种气体里烧得格外旺。他给这种气体起名"氧"。结论干净利落：燃烧不是放出燃素，而是物质\textbf{与氧结合}，增加的质量来自空气。密封容器里反应前后总质量不变——这正是第 26 章质量守恒定律最经典的实验来源之一。

值得多说一句的是，这种气体拉瓦锡并不是第一个制出来的人。英国的普利斯特里和瑞典的舍勒都更早得到了氧气，普利斯特里甚至发现蜡烛在这种气体里烧得特别亮、老鼠在里面活得特别精神——但他用燃素说的框架去解释，管它叫"脱燃素空气"：不是它帮助了燃烧，而是它"特别能吸收燃素"。同样的实验现象，装进不同的理论框架，就讲出完全不同的故事。这个对比是科学史上最好的课之一：\textbf{实验数据不自带解释}，看到什么很重要，用什么概念去整理看到的，同样重要。

故事到这里只讲了一半。拉瓦锡回答了"质量从哪来"，却没回答"结合的本质是什么"——他不可能回答，因为电子还要再等一百多年（1897 年）才被发现。20 世纪，当电子转移的图景建立起来，化学家回头一看：氧之所以在燃烧中唱主角，只是因为它抢电子的本领高；燃烧、生锈、呼吸的共同本质是电子转移。于是"氧化"这个词从"和氧结合"被\textbf{重新定义}为"失去电子"。科学概念就是这样：词留下来了，含义换掉了。这也提醒我们，学一个科学名词，要追的是它今天指什么，而不是它字面上像什么。
\end{zhengju}

\begin{shiyan}
\textbf{铁钉为什么会锈：三组对照}

材料：三枚干净无锈的铁钉（用砂纸打磨光亮）、三个透明杯子、自来水、食盐、食用油、干燥剂（零食袋里的小包即可）、保鲜膜。

步骤：① 杯 1 装半杯自来水，放入铁钉，敞口放置——钉子一半泡在水里一半接触空气。② 杯 2 装同样多的水，加一勺盐搅溶，放入铁钉。③ 杯 3 不放水，放铁钉和两包干燥剂，杯口用保鲜膜封紧。如果家里有食用油，可以再加一杯 4：水面上倒一层油封住水面再放钉子，试试"有水没空气"会怎样。

观察点：三天内每天看一次，记录哪个杯子里先出现红棕色锈斑、锈长在哪里（水面线附近？还是杯底？）。

预期与思考：杯 2 通常锈得最快，杯 3 几乎不锈。这说明铁生锈\textbf{同时需要水和氧气}，盐会大大加速它——电子在盐水里"跑路"更顺畅。想一想：为什么杯 1 里水面线附近锈得最重？（提示：那里氧气和水同时最充足。）

安全提示：实验本身安全，但用过的铁钉边缘可能锋利，处理时垫纸拿取；打完砂纸的碎屑及时清理，别弄到眼睛里。
\end{shiyan}

\section{这套记账本在哪里会算错账}

氧化数好用，但请记住它"假装清算"的本质，三种情况下它会露出马脚。

第一，它给出的常常不是真实电荷。\chem{CO_2} 里的碳记作 $+4$，但它并不带 $+4$ 的电荷——真实的电荷偏移要小得多，而且连续分布在原子和键上。拿氧化数去算真实的静电作用力，会算出笑话。

第二，有些分子里平均氧化数会是分数。四氧化三铁 \chem{Fe_3O_4} 里铁的平均氧化数是 $+8/3$。没有"三分之八价"的原子，真相是这个晶体里住着两种铁：一部分 $+2$、一部分 $+3$，平均成了分数。账本给的是平均数，货架上其实有两种商品。

第三，判断"是不是氧化还原"要逐元素看账，不能凭感觉。碳酸钙分解成氧化钙和二氧化碳，看起来"大变样"，可每个元素的氧化数都没动——没有电子转移，它不是氧化还原。反过来，酸碱中和（第 32 章）也不是氧化还原：质子转移了，电子没转移。化学把反应分成"电子转移推动的"和"不转移电子的"两大类，这个分法比"有没有氧气""有没有火焰"可靠得多。

\begin{wuqu}
"氧化反应就是物质和氧气发生的反应。"——这是"氧化"这个倒霉名字制造的最大误解。反例随手就有：氢气在氯气中燃烧，生成氯化氢，全程没有氧，氢照样被氧化（氧化数 0 升到 $+1$）；铜丝插进硝酸银溶液，铜被氧化成 \chem{Cu^{2+}}，氧根本没露面；你体内的铁离子在 $+2$ 和 $+3$ 之间转换，也不必有氧气直接参与。判断的唯一标准是电子（或氧化数）：失电子、氧化数升高，就是被氧化。氧气只是众多"电子买家"里最出名的一个。顺带一提，它的反向误解也存在："还原"不是"变回原来的东西"，而是\textbf{得到电子}——铁矿石被还原成铁，得到电子的是铁离子，"还"掉的其实是氧。
\end{wuqu}

\section{回到那枚铁钉}

现在回答开场的三个问题。

第一，铁锈比原来的铁\textbf{更重}——多出来的是从空气里结合的氧。铁钉整体看起来变轻，只是因为锈壳疏松剥落、掉渣了；如果把掉下的渣都收集起来称，账是平的，而且是增加的。拉瓦锡的天平早在两百多年前就称出了这个结论。

第二，红棕色的壳不是铁"变质"，也不是外面贴的灰。它是铁原子失去电子变成铁离子、再与氧和水结合生成的水合氧化铁。铁元素还在，只是电子状态变了，邻居换了——就像第 10 章说的，同一张"身份证"，不同的角色。

第三，苹果和铁钉，是同一出戏。苹果切面里的酚类物质遇到空气中的氧，在细胞释放出的酶（第 49 章）的帮助下被氧化，生成黄褐色的产物——你看到的是一场缓慢的、被酶加速的氧化。泡盐水、滴柠檬汁之所以管用，是因为它们干扰了酶的工作或改变了环境酸度，让这场氧化减速。铁生锈也是氧化，但没有酶帮忙，慢得本分。燃烧是快到发光的氧化，生锈和苹果变褐是慢到几乎看不见的氧化，而你的呼吸——下面就要讲——是被生命\textbf{精确控制了速度}的氧化。

\section{从铁锈到呼吸：电子转移无处不在}

把镜头拉远，你会发现电子转移是这颗星球上最繁忙的生意之一。

冶炼厂干的活，本质上是用强还原剂把电子"塞回"给矿石里的金属离子：高炉里，一氧化碳把电子交给氧化铁里的铁离子，铁被还原成铁水流出（第 15 章的铁，就是这样从矿石里"赎"出来的）。漂白剂和消毒剂干的是相反方向的活：次氯酸钠是凶猛的氧化剂，把色素分子和微生物的关键分子强行氧化、破坏结构，颜色褪去，细菌丧命（所以漂白剂绝不能和酸性洁厕灵混用——会放出氯气，这是氧化还原在家里最危险的用法）。

厨房里天天上演的燃烧，也可以把账记一遍。天然气的主要成分甲烷 \chem{CH_4} 燃烧：\[\chem{CH_4} + \chem{2O_2} \rightarrow \chem{CO_2} + \chem{2H_2O}\]记一记氧化数：甲烷里的碳是 $-4$（氢 $+1$、电负性更大的碳把共享电子判给自己），二氧化碳里的碳是 $+4$——一个碳原子交出 8 个电子；氧从 0 降到 $-2$，两个氧分子共收下 8 个电子，账正好平。蓝色的火苗，就是这笔每秒亿万次成交的电子交易所发的光。燃气灶要留通风口、燃烧不充分会产生一氧化碳，这些安全常识的底层，都是这同一笔账：没有足够的氧化剂，碳只能"交一半"电子，停在有毒的中间产物上。

而最重要的那场氧化还原，此刻正在你身体里进行。你吃下的葡萄糖，碳的平均氧化数是 0；你呼出的二氧化碳里，碳是 $+4$。从嘴到肺，每个碳原子交出了 4 个电子——和燃烧葡萄糖\textbf{完全同一笔总账}。区别只在于：篝火让电子"一拥而上"，能量以光和热的形式浪费掉；细胞则修了一条几十步长的"电子传送链"，让电子一小步一小步地走，每一步释放的一小口能量都被收集起来，存进 ATP。呼吸不是"吸进氧气吐出二氧化碳"这么简单，氧气只是站在传送链终点、最后收下电子并生成水的那个"终点站"（第 52、53 章会把这条链拆开给你看）。光合作用则把整笔账反过来做：用阳光当动力，把电子硬塞回二氧化碳，重新造成糖——而氧气来自被拆开的水分子，不是二氧化碳（第 54 章，这条是全书红线，先记住结论）。

你看，从铁钉到高炉，从苹果到肺，氧化还原是一条贯穿无机界和生命界的主线。而这条主线上还有一个诱人的岔口，本章已经埋过一次：半反应既然能分开写，就能在空间上分开做——把氧化和还原分别关在两个"房间"里，逼电子走一条你指定的外电路，电子在路上就能替你做功：点亮灯泡、驱动马达、给手机充电。化学变化就这样被改造成了电流。下一章，我们就去造一座电池。

\begin{tiaozhan}
试试用半反应法配平一个经典反应：紫红色的高锰酸钾溶液遇到草酸 \chem{H_2C_2O_4}（菠菜、茶叶里都有）会褪色。先写两个半反应。锰这边：\chem{MnO_4^-} 变成 \chem{Mn^{2+}}，氧化数从 $+7$ 降到 $+2$，收 5 个电子；左边氧多，用 \chem{H^+} 和水配平原子：\[\chem{MnO_4^-} + \chem{8H^+} + \chem{5e^-} \rightarrow \chem{Mn^{2+}} + \chem{4H_2O}\]碳这边：草酸里碳是 $+3$，变成 \chem{CO_2} 里的 $+4$，两个碳共交 2 个电子：\[\chem{H_2C_2O_4} \rightarrow \chem{2CO_2} + \chem{2H^+} + \chem{2e^-}\]最小公倍数是 10：第一式乘 2、第二式乘 5，相加消去电子，整理即得配平结果。数一数两边每种原子和总电荷，账应该是平的。还有一个有趣的现象：一个分子里同一元素一部分升、一部分降的反应叫\textbf{歧化反应}，比如过氧化氢分解（氧从 $-1$ 分别变成 $-2$ 和 0）——它同时是自己的氧化剂和还原剂。跳过本节不影响主线。
\end{tiaozhan}

\section*{练一练}

\begin{sikaoti}
\item 查账判断：下列过程哪些是氧化还原？① 水结成冰；② 天然气（甲烷 \chem{CH_4}）燃烧；③ 盐酸和氢氧化钠中和；④ 绿色植物的光合作用。对是的那些，指出什么元素被氧化、什么被还原。
\item 画出"铜丝插入硝酸银溶液"的粒子示意图：标出铜原子、铜离子、银离子、银原子，用箭头画出电子走的路线，并在图旁写明两个半反应。注明：圆球大小、颜色都是人为的表示法。
\item 用氧化数记账：\chem{SO_2}、\chem{SO_4^{2-}}、\chem{H_2S} 中硫的氧化数各是多少？火山口的硫从 \chem{H_2S} 被氧化成 \chem{SO_2}，每个硫原子交出了几个电子？
\item 家里同时有"84 消毒液"（次氯酸钠，氧化剂）和"洁厕灵"（盐酸）。查阅它们的成分后解释：为什么混合会产生氯气？从氧化还原角度说说氯在这里经历了什么。（只查资料，\textbf{严禁}实际混合验证。）
\item 画一幅"电子流图"：把葡萄糖燃烧（篝火）和葡萄糖在细胞内被分解（呼吸）画成两条路径，起点和终点相同，路径不同。在图上标出：哪条路径的能量是一下子放掉的，哪条被分成了许多小步。思考：同样的总账，为什么细胞的方式能"攒下钱"？
\item 有人说"生锈是铁被氧化了，所以铁锈是铁'变坏'了，想办法把氧挡在外面就永远不用管"。请用本章的模型评论这句话：前半句对不对？后半句的思路在工程上可行吗（想想轮船、桥梁、食品包装各用什么办法）？还有没有"让铁生锈也没关系"的情况？
\end{sikaoti}
